% 填入你的論文題目、姓名等資料
% 如果題目內有必須以數學模式表示的符號，請用 \mbox{} 包住數學模式，如下範例
% 如果中文名字是單名，與姓氏之間建議以全形空白填入，如下範例
% 英文名字中的稱謂，如 Prof. 以及 Dr.，其句點之後請以不斷行空白~代替一般空白，如下範例
% 如果你的指導教授沒有如預設的三位這麼多，則請把相對應的多餘教授的中文、英文名
%    的定義以空的大括號表示
%    如，\renewcommand\advisorCnameB{}
%          \renewcommand\advisorEnameB{}
%          \renewcommand\advisorCnameC{}
%          \renewcommand\advisorEnameC{}


% 論文題目 (中文)
\renewcommand\cTitle{
可驗證之安全系統的應用
}

% 論文題目 (英文)
\renewcommand\eTitle{
A Verifiable Secure protocol in a Secure System 
}

% 我的姓名 (中文)
\renewcommand\myCname{王大明}

% 我的姓名 (英文)
\renewcommand\myEname{Da-Ming Wang}

% 我的學號
\renewcommand\myStudentID{M19150001}

% 指導教授A的姓名 (中文)
\renewcommand\advisorCnameA{陳明明\ 博士}

% 指導教授A的姓名 (英文)
\renewcommand\advisorEnameA{Ming-Ming Cheng, Ph.D.}

% 指導教授B的姓名 (中文)
\renewcommand\advisorCnameB{}

% 指導教授B的姓名 (英文)
\renewcommand\advisorEnameB{}

% 指導教授C的姓名 (中文)
\renewcommand\advisorCnameC{}

% 指導教授C的姓名 (英文)
\renewcommand\advisorEnameC{}

% 校名 (中文)
\renewcommand\univCname{國立雲林科技大學資訊工程系}

% 科系 (英文)
\renewcommand\univdepartmentEname{Department of Computer Science and Information Engineering}

% 校名 (英文)
\renewcommand\univEname{National Yunlin University of Science \& Technology \\
Master Thesis
}

% 學位名 (中文)
\renewcommand\degreeCname{碩士}

% 口試年份 (中文、民國)
\renewcommand\cYear{115}

% 口試月份 (中文)
\renewcommand\cMonth{6} 

% 口試年份 (阿拉伯數字、西元)
\renewcommand\eYear{2026} 

% 口試月份 (英文)
\renewcommand\eMonth{June}
